\chapter{Background}
\label{ch:background}

This chapter establishes the theoretical foundation for the dissertation. It is organized into two main parts. Part I provides the conceptual background on machine learning, workload prediction, and intelligent systems. Part II reviews the specific state-of-the-art literature and identifies the research gaps addressed in Chapters~\ref{ch:adaptive-orchestration} through \ref{ch:distributed-agentic-ai}.

A unifying thread connects these foundations: the attention mechanism. Originally developed for sequence-to-sequence tasks, attention has become the architectural core of transformer-based models used in time-series forecasting (Chapter~\ref{ch:attention-driven-prediction}) and the dominant \acp{LLM} that power autonomous agents (Chapter~\ref{ch:distributed-agentic-ai}). By establishing attention fundamentals in the context of workload prediction, this chapter provides the conceptual basis for understanding both the prediction and autonomy components of this dissertation.

% Part I: Foundational Concepts
\section{Part I: Foundational Concepts}
\label{sec:ch2_part1}

This section introduces the core technical concepts required to understand the proposed solutions: workload prediction, generalization mechanisms, and intelligent orchestration architectures.

% 2.1 Workload Prediction and Resource Management
\subsection{Workload Prediction and Resource Management}
\label{sec:ch2_workload_prediction}

This section examines workload prediction and resource management in edge-cloud environments, first examining how workload predictions drive resource management decisions, then distinguishing between \ac{VM}-level and container-level prediction requirements. Throughout, we identify specific limitations that motivate the contributions in Chapters~\ref{ch:adaptive-orchestration} and \ref{ch:attention-driven-prediction}.

\subsubsection{Resource Management and Orchestration}
\label{sec:ch2_resource_mgmt}

Workload predictions drive proactive resource management decisions across the edge-cloud continuum.

Predicted workload increases trigger three orchestration actions: horizontal scaling (adding container replicas), vertical scaling (increasing resource allocations), or migration (relocating containers). Horizontal scaling deploys additional container replicas when:

\begin{equation}
\hat{\mathbf{u}}_i(t+\delta) > \tau_s \cdot \mathbf{r}_i,
\end{equation}

where $\mathbf{r}_i$ is node capacity, $\hat{\mathbf{u}}_i(t+\delta)$ is predicted utilization, and $\tau_s$ is the scale-out threshold. Production clusters typically maintain machine utilization below 80\% to preserve scheduling headroom~\cite{tirmazi2020borg}. Vertical scaling increases \ac{CPU}/memory allocations when average utilization exceeds thresholds, avoiding container restarts when possible. Migration relocates containers to balance load or consolidate resources for energy efficiency, with typical migration time $t_m \in [8, 30]$ seconds~\cite{eskandani2024criu}. These orchestration actions can be triggered either reactively or proactively.

Reactive orchestration~\cite{8258257} responds to current metrics, often triggering after \ac{QoS} violations occur. Proactive orchestration~\cite{farhad2023ai} uses predictions to make decisions before violations occur:

\begin{equation}
t_d = t_c + \delta - t_a,
\end{equation}

where $t_d$ is the decision time, $t_c$ is current time, $\delta$ is the prediction horizon, and $t_a$ is the action execution time. Accurate prediction with horizon $\delta \geq t_a$ enables zero-downtime elasticity.

While proactive orchestration enables timely decisions, these decisions must also balance competing objectives. Resource management involves critical trade-offs quantified by the cost function:

\begin{equation}
J_{total} = \lambda J_{qos} + \mu E(\pi, t) + \gamma R(\pi, t),
\end{equation}

where $J_{qos}$ represents \ac{QoS} penalty, $E(\pi, t)$ is energy cost, $R(\pi, t)$ is load balancing cost, and $\lambda$, $\mu$, $\gamma$ are trade-off weights. Conservative predictions minimize \ac{SLA} violations but increase over-provisioning costs (industry studies show 30-50\% rates~\cite{raith2024opportunistic}). Aggressive predictions reduce over-provisioning but risk expensive \ac{SLA} violations. Effective resource management requires balancing prediction accuracy and computational efficiency for real-time edge decisions. These requirements manifest differently depending on whether workload forecasting operates at the \ac{VM} or container level.

\subsubsection{Container-Level vs. \acs{VM}-Level Prediction}
\label{sec:ch2_container_prediction}

Traditional cloud computing relies on \ac{VM} abstraction, while modern edge-cloud systems use containerized microservices. This architectural shift introduces unique prediction challenges.

\ac{VM}-level characteristics~\cite{tirmazi2020borg} include: (i) long-lived instances (hours to weeks), (ii) coarse-grained resource allocation (whole \ac{CPU} cores, \ac{GB} memory), (iii) lower deployment density (tens per physical host), (iv) rich historical data for training, and (v) relatively stable resource consumption patterns. In contrast, container-level characteristics~\cite{10202641, 10417087, 10029931} include: (i) short-lived instances (minutes to hours), (ii) fine-grained resource sharing (fractional \ac{CPU}, \ac{MB} memory), (iii) high deployment density (hundreds per host), (iv) sparse historical data per container, and (v) highly variable patterns due to microservice dependencies.

This data scarcity can be quantified as:

\begin{equation}
N_c = \frac{L_c}{s} \ll \frac{L_{vm}}{s} = N_{vm},
\end{equation}

where $L_c$ and $L_{vm}$ are container and \ac{VM} lifetimes, $s$ is the sampling interval, and $N_c$, $N_{vm}$ are available training samples. Container lifetimes are substantially shorter than \acp{VM}~\cite{nadgowda2017voyager}, providing insufficient training data for traditional deep learning models.

Container-level prediction requires models that: (i) can learn effectively from limited data, (ii) generalize across diverse container types, (iii) handle high-frequency sampling rates (seconds vs. minutes), and (iv) account for inter-container dependencies in microservice architectures. Traditional \ac{VM}-level predictors struggle in container environments because they assume data abundance and pattern stability.

Chapters~\ref{ch:adaptive-orchestration} and \ref{ch:attention-driven-prediction} specifically target container-level workload prediction, developing lightweight models capable of learning from sparse traces and generalizing across heterogeneous container deployments through architectural innovations rather than increased data requirements. These innovations build upon machine learning architectures suited to time-series forecasting under resource constraints.

\subsubsection{Machine Learning for Workload Prediction}
\label{sec:ch2_ml_workload}

Traditional rule-based orchestration systems struggle with dynamic edge-cloud networks where workload patterns evolve continuously~\cite{chen2018auto}. Machine learning addresses these limitations by learning patterns from historical data, enabling predictive and autonomous decision-making without manual feature engineering.

For workload prediction, let $\mathbf{W}_t = [W_t^{(1)}, W_t^{(2)}, \ldots, W_t^{(f)}]^\top \in \mathbb{R}^f$ represent the workload vector at time step $t$, where $W_t^{(i)}$ denotes the $i$-th resource feature (\ac{CPU}, memory usage). The goal is to predict future workload vectors based on historical observations using a prediction function:

\begin{equation}
\hat{\mathbf{W}}_{t+1:t+\lambda} = \Psi_{\boldsymbol{\phi}}\left( \mathbf{W}_{t-\kappa+1:t} \right),
\end{equation}

where $\kappa$ is the lookback window, $\lambda$ is the prediction horizon, and $\boldsymbol{\phi}$ represents learnable model parameters.

Modern machine learning architectures address this through different approaches. \textit{Sequential models} (\acp{RNN}, \ac{LSTM}) maintain hidden states across time steps but suffer from parallelization limitations~\cite{hochreiter1998vanishing}. \textit{Convolutional models} (\acp{TCN}) apply learned filters across time windows with efficient parallel computation but limited receptive fields~\cite{donghao2024moderntcn}. \textit{Attention-based models} (Transformers) dynamically focus on relevant portions of input sequences but introduce quadratic complexity~\cite{vaswani2017attention}. \textit{Feed-forward models} (\acp{MLP}) offer minimal computational overhead suitable for edge deployment but assume fixed periodicities~\cite{sparseTSF}.

Models are trained by minimizing prediction loss, typically Mean Squared Error:

\begin{equation}
\mathcal{L}(\boldsymbol{\phi}) = \frac{1}{\lambda} \sum_{i=1}^{\lambda} \left\| \mathbf{W}_{t+i} - \hat{\mathbf{W}}_{t+i} \right\|_2^2,
\end{equation}

Training uses gradient-based optimizers to minimize this loss. However, standard deep learning approaches conflict with edge requirements. For edge deployment, critical constraints include model size (typically <1\ac{GB} memory), inference latency (<50ms for real-time orchestration~\cite{8334540, 9500858}), and accuracy sufficient for proactive resource management. These constraints exclude most deep learning architectures from edge consideration, motivating the lightweight adaptive approaches developed in Chapters~\ref{ch:adaptive-orchestration} and \ref{ch:attention-driven-prediction}.

Beyond computational constraints, production deployments face an additional challenge: models trained on one workload type often fail when applied to different applications, creating a generalization barrier that compounds the data scarcity problem identified in Section~\ref{sec:ch2_container_prediction}.
% 2.2 Generalization
\subsection{Generalization}
\label{sec:ch2_generalization}

The container-level prediction challenges identified in Section~\ref{sec:ch2_workload_prediction}, particularly the sparse training data and high workload variability, are compounded when models must generalize across different application types. This section introduces the challenge of model generalization in edge-cloud environments, explains the attention mechanism for capturing temporal dependencies, and presents temporal clustering as a method for handling workload heterogeneity. These concepts establish the theoretical basis for the \ac{OmniFORE} framework presented in Chapter~\ref{ch:attention-driven-prediction}.

\subsubsection{The Cross-Dataset Generalization Challenge}
\label{subsec:ch2_generalization_challenge}

Models trained on one workload dataset often show degraded performance when deployed on different workload types, even within the same infrastructure~\cite{zou2024generalization}. This occurs because edge-cloud workloads show fundamental heterogeneity across two dimensions. The first dimension is \textit{workload diversity}: web services, databases, batch processing, and streaming applications generate distinct resource consumption patterns. These patterns also vary across edge, fog, and cloud tiers due to hardware and network differences. The second dimension is \textit{temporal variability}. Resource usage at time $t$ can be modeled as
\[
T_i(t) = \hat{T}_i(t) + \epsilon,
\]
where $\hat{T}_i(t)$ captures the predictable component and $\epsilon$ represents stochastic variations from user demand fluctuations, scheduled tasks, and time-of-day effects. This heterogeneity violates the i.i.d. assumption of traditional machine learning.

Formally, consider a model $f_\theta$ trained on source dataset $\mathcal{D}_{\text{source}}$ and evaluated on target dataset $\mathcal{D}_{\text{target}}$. The generalization gap is:
\begin{equation}
\Delta = \mathcal{L}(f_\theta; \mathcal{D}_{\text{target}}) - \mathcal{L}(f_\theta; \mathcal{D}_{\text{source}}),
\end{equation}
where $\mathcal{L}$ is the prediction loss. Large positive $\Delta$ indicates poor generalization, requiring retraining or adaptation for each new deployment.

This generalization gap creates operational challenges. Consider a production edge-cloud network serving diverse application types: web frontends, \ac{API} gateways, databases, caches, batch processors, and stream processors. Without generalization, practitioners must maintain separate models for each application category. Each model requires independent training, separate monitoring configurations, and coordinated updates to prevent version conflicts. When a new application type deploys, operators face a choice: deploy without prediction (reverting to reactive orchestration with its associated over-provisioning costs), or delay deployment until sufficient training data accumulates and a new model trains. This operational overhead makes prediction-based orchestration more challenging at scale.

Addressing this generalization gap requires prediction approaches that capture common temporal patterns across diverse workloads rather than memorizing dataset-specific characteristics. Recent advancements in machine learning for time-series forecasting have introduced techniques that can identify complex dependencies across different timescales. Attention mechanisms have shown effectiveness in generalization across diverse time-series patterns~\cite{vaswani2017attention, wen2022transformers}. Notably, the same attention architecture that enables workload generalization also powers dominant \acp{LLM}~\cite{vaswani2017attention}, creating a conceptual bridge between the prediction techniques in this dissertation and the autonomous orchestration systems that consume these predictions. Attention-enhanced neural networks and attention-based transfer learning have shown particular promise for workload prediction, selectively focusing on critical temporal patterns while filtering out noise. Understanding how attention mechanisms achieve this selective focus requires examining their mathematical foundations.

\subsubsection{Introduction to Attention Mechanism}
\label{subsec:ch2_attention_mechanism}

The attention mechanism, introduced by Vaswani et al.~\cite{vaswani2017attention}, represents a foundational advance in sequence modeling. Beyond its role in time-series forecasting, attention is the architectural core of dominant \acp{LLM}, including the GPT family, that power the autonomous agents discussed in Chapter~\ref{ch:distributed-agentic-ai}. This connection arises from a shared abstraction: both time-series data and natural language are sequences of vectors. A workload trace $\mathbf{W}_{t-\kappa:t}$ and a text prompt are both ordered collections of embeddings that attention processes using the same underlying mechanism. By establishing attention fundamentals in the time-series context, this section provides the conceptual foundation for understanding both workload prediction (Chapter~\ref{ch:attention-driven-prediction}) and \ac{LLM}-based orchestration (Chapter~\ref{ch:distributed-agentic-ai}).

The attention mechanism addresses the generalization challenge by learning which temporal patterns are relevant for prediction rather than relying on fixed feature extractors. This data-driven relevance weighting enables models to identify similar patterns across diverse workloads, even when those workloads differ in absolute magnitude or baseline characteristics.

Before attention can operate, raw input must be converted into embeddings. For a time-series, each time step's feature vector $\mathbf{x}_t \in \mathbb{R}^{d_f}$ (e.g., CPU and memory utilization) is projected into a higher-dimensional embedding space via a learned linear transformation: $\mathbf{E}_t = \mathbf{x}_t \mathbf{W}_E + \mathbf{p}_t$, where $\mathbf{W}_E \in \mathbb{R}^{d_f \times d_{\text{embed}}}$ is a weight matrix and $\mathbf{p}_t \in \mathbb{R}^{d_{\text{embed}}}$ is a positional encoding that preserves temporal order~\cite{vaswani2017attention}. For text sequences, the process is analogous: each token is mapped to a vector through an embedding lookup table, then combined with positional information. In both cases, the result is a sequence of embeddings $\mathbf{E} \in \mathbb{R}^{d_x \times d_{\text{embed}}}$ that serves as input to the attention mechanism.

Starting with these initial embeddings, where $d_x$ is the input sequence length and $d_{\text{embed}}$ is the embedding dimension, the attention mechanism (denoted as $\mathcal{A}$) refines them into a new set, $\mathbf{E}'$~\cite{vaswani2017attention}. This mechanism allows the model to selectively focus on specific time steps, similar to human attention. For instance, when predicting future workload at a peak, both the model and a human analyst would focus on the pattern surrounding the current peak and compare it to similar peaks in the past. The attention weights determine how much each original embedding $\mathbf{E}_i \in \mathbb{R}^{d_{\text{embed}}}$ contributes to its corresponding refined embedding $\mathbf{E}_i'$. As shown in Figure~\ref{fig:attention_goal}, thicker lines indicate stronger influence. The refined embeddings $\mathbf{E}'$ effectively combine underlying data patterns with temporal relationships, improving the model's ability to capture complex dependencies across the sequence~\cite{dai2019transformer}.

\begin{figure}[!htbp]
\centering
\includegraphics[width=0.866\textwidth]{sections/chapter2/img/attention_goal.pdf}
\caption{Attention mechanism refining time-series embeddings $\mathbf{E}_i$ to $\mathbf{E}_i'$, with arrow thickness indicating the strength of influence.}
\label{fig:attention_goal}
\end{figure}

To construct the refined embeddings, the attention mechanism uses three key components: queries, keys, and values~\cite{vaswani2017attention}, as illustrated in Figure~\ref{fig:attention_workings}. For each row $\mathbf{E}_i \in \mathbb{R}^{d_{\text{embed}}}$ of the embedding matrix $\mathbf{E}$ (corresponding to time step $i$), the model computes three vectors:

\begin{equation}
    \begin{aligned}
        \text{Query vector:} \quad \mathbf{Q}_i &= \mathbf{E}_i \mathbf{W_Q}, \\
        \text{Key vector:} \quad \mathbf{K}_i &= \mathbf{E}_i \mathbf{W_K}, \\
        \text{Value vector:} \quad \mathbf{V}_i &= \mathbf{E}_i \mathbf{W_V},
    \end{aligned}
\end{equation}
where $\mathbf{W_Q}, \mathbf{W_K}, \mathbf{W_V} \in \mathbb{R}^{d_{\text{embed}} \times d_{\text{embed}}}$ are learned weight matrices. Conceptually, a query vector encodes the information need at the current time step, key vectors help identify relevant past information, and value vectors contain the actual content to be aggregated~\cite{bahdanau2014neural}. In the context of workload prediction, when analyzing a peak, the query vector would encode features of the current peak, key vectors would help identify similar historical peaks, and value vectors would contain the associated historical workload patterns.

The model evaluates temporal relationships by computing scaled dot products between query and key vectors:

\begin{equation}
    \text{Score}(\mathbf{Q}_i, \mathbf{K}_j) = \frac{\mathbf{Q}_i \mathbf{K}_j^T}{\sqrt{d_{\text{embed}}}}.
\end{equation}

A high score indicates that the corresponding time step contains critical information for the current prediction, as the vectors align in high-dimensional space~\cite{luong2015effective}. The scaling factor $\sqrt{d_{\text{embed}}}$ ensures training stability by preventing excessively large values that could lead to sharp, unbalanced distributions post-softmax~\cite{xu2015show}. After scaling, the softmax function normalizes these scores into a probability distribution, effectively assigning weights to each timestep's embedding vector. These weights, summing to one, determine the relative importance of each timestep, guiding the attention mechanism to focus on the most relevant parts of the sequence.

To capture diverse temporal dependencies, the model uses \textit{multi-headed attention}~\cite{michel2019sixteen} with $H$ parallel attention mechanisms, or "heads" (Fig.~\ref{fig:attention_workings}). Each head $h \in \{1, ..., H\}$ operates independently on the same input embeddings, but with its own set of learnable weight matrices $\mathbf{W}_Q^h, \mathbf{W}_K^h, \mathbf{W}_V^h \in \mathbb{R}^{d_{\text{embed}} \times d_{\text{head}}}$, where $d_{\text{head}} = d_{\text{embed}} / H$. These matrices are initialized differently, leading to distinct evolutionary trajectories during training and convergence to different local optima~\cite{zhang2019improving}. This architectural design allows the model to simultaneously project the input into multiple subspaces and attend to diverse aspects of the sequence, such as short-term fluctuations and long-term trends. This effectively increases the model's representational capacity without a proportional increase in computational complexity~\cite{vaswani2017attention}.

For each head $h$, the attention output is computed as

\begin{equation}
    \mathbf{A}_i^h = \mathcal{A}(\mathbf{E}_i\mathbf{W}_Q^h, \mathbf{E}_i\mathbf{W}_K^h, \mathbf{E}_i\mathbf{W}_V^h).
\end{equation}

The outputs from all heads are concatenated and linearly transformed to produce the \ac{MHA} output

\begin{equation}
    \mathbf{MHA}_i = \operatorname{Concat}(\mathbf{A}_i^{(1)}, \mathbf{A}_i^{(2)}, \ldots, \mathbf{A}_i^{(H)}) \mathbf{W}^O,
\end{equation}
where $\mathbf{W}^O \in \mathbb{R}^{d_{\text{embed}} \times d_{\text{embed}}}$ is a learned output projection matrix. Since each head outputs a vector of dimension $d_{\text{head}} = d_{\text{embed}}/H$, concatenating all $H$ heads yields a vector of dimension $H \times d_{\text{head}} = d_{\text{embed}}$, which is then projected back to $d_{\text{embed}}$.

The final refined embedding for each time step incorporates a residual connection~\cite{he2016deep} and layer normalization~\cite{lei2016layer}:

\begin{equation}
    \mathbf{E}_i' = \text{LayerNorm}(\mathbf{E}_i + \mathbf{MHA}_i),
\end{equation}
where \ac{LayerNorm} denotes layer normalization. This is followed by a position-wise feed-forward network ($d_{ff}$) with another residual connection and layer normalization:

\begin{equation}
    \mathbf{E}_i'' = \text{LayerNorm}(\mathbf{E}_i' + d_{ff}(\mathbf{E}_i')).
\end{equation}

This entire process, comprising multi-headed attention, residual connections, layer normalizations, and the feed-forward network, constitutes one layer of the transformer encoder. It is repeated across multiple stacked layers ($N_e$). Each layer operates on the output of the previous one, allowing for iterative refinement of embeddings and progressively developing detailed representations needed for accurate predictions in complex time-series tasks.

\begin{figure}[!htbp]
\centering
\includegraphics[width=0.595\textwidth]{sections/chapter4/img/attention_workings.pdf}
\caption{Multi-Headed Attention: Input embeddings ($\mathbf{E}_i$) are transformed into query ($\mathbf{Q}$), key ($\mathbf{K}$), and value ($\mathbf{V}$) vectors. Multiple heads ($H$) capture different temporal aspects. Larger circles indicate stronger attention scores.}
\label{fig:attention_workings}
\end{figure}

The attention mechanism provides the mathematical framework for capturing temporal dependencies across sequences. Applying this framework effectively to the diverse workloads characterized in Section~\ref{subsec:ch2_generalization_challenge} requires a method for organizing training data that covers this heterogeneity.

\subsubsection{Temporal Clustering for Workload Heterogeneity}
\label{subsec:ch2_temporal_clustering}

Attention mechanisms enable models to learn relevant temporal patterns, but training such models still requires diverse workload examples. Rather than requiring exhaustive data from every possible workload type, temporal clustering can identify representative workload categories. This reduces training requirements while maintaining coverage of diverse patterns.

However, direct clustering of high-dimensional time-series data is computationally expensive and often struggles to capture meaningful patterns. For a set of container traces $\mathbf{T}(t) = \{T_1, T_2, \ldots, T_{N_{\text{all}}}\}$ where each trace $T_i(t) \in \mathbb{R}^{d_f \times d_t}$, the raw dimensionality $d_f \times d_t$ makes direct clustering inefficient. Latent space transformation addresses this challenge by projecting traces into lower-dimensional representations that preserve essential temporal patterns:
\[
\mathcal{Z}(z) = f_{\text{encoder}}(\mathcal{T}(t)), \quad \text{where } z \ll d_f \times d_t.
\]

The encoder function $f_{\text{encoder}}$ compresses the high-dimensional time-series into a dense latent representation $\mathcal{Z}$, enabling efficient clustering into $K$ distinct groups:
\[
\{C_1, C_2, \ldots, C_K\} = \text{Cluster}(\mathcal{Z}(z)).
\]

Each cluster $C_k$ groups container traces with similar spatio-temporal patterns, addressing workload heterogeneity by identifying fundamental pattern categories that generalize across applications. This approach allows models to learn from representative samples of each cluster rather than needing exhaustive training on all individual traces, improving computational efficiency while maintaining generalization capability.

This clustering-based approach directly addresses the cross-dataset generalization challenge identified in Section~\ref{subsec:ch2_generalization_challenge}. By grouping workloads into representative clusters and training attention-based models on diverse samples from each cluster, systems can generalize across applications without requiring exhaustive per-application training. Chapter~\ref{ch:attention-driven-prediction} implements these principles through the \ac{OmniFORE} framework, combining temporal clustering with attention mechanisms for cross-application workload prediction.

With workload prediction (Section~\ref{sec:ch2_workload_prediction}) and generalization mechanisms established, the remaining foundational challenge is translating predictions into orchestration actions that satisfy operator intent.

% 2.3 Intelligent Orchestration Systems
\subsection{Intelligent Orchestration Systems}
\label{sec:ch2_intelligent_orchestration}

While Section~\ref{sec:ch2_workload_prediction} established how predictions enable proactive resource management, and Section~\ref{sec:ch2_generalization} addressed scaling these predictions across diverse workloads, a critical gap remains in translating high-level operator intent into concrete orchestration actions. This section establishes mathematical foundations for AI-driven orchestration, covering the end-to-end orchestration problem and its three stages, placement as constrained optimization, \ac{LLM}-based reasoning with structured outputs, autonomous agent formalization, and multi-agent coordination mechanisms essential for understanding Chapter~\ref{ch:distributed-agentic-ai}.

\subsubsection{The End-to-End Orchestration Problem}

Real-world orchestration extends beyond optimization to encompass three distinct stages. In Stage~(i), \textit{intent interpretation}, natural language requests (e.g., ``reduce energy while maintaining \ac{QoS}'') must be translated into formal objectives and constraints. In Stage~(ii), \textit{placement optimization}, a placement policy is derived that satisfies those constraints. In Stage~(iii), \textit{compliance monitoring}, infrastructure is continuously observed for unexpected state changes and intent violations. Traditional optimization methods such as \ac{ILP}~\cite{margot2009symmetry}, genetic algorithms~\cite{holland1992genetic}, and deep reinforcement learning~\cite{li2017deep} primarily address Stage~(ii), assuming objectives and constraints are already formalized and states are static. This gap motivates AI-driven approaches that can handle the full pipeline. Among these three stages, placement optimization (Stage~(ii)) represents the mathematical core that any orchestration approach must address.

\subsubsection{Placement as Constrained Optimization}

Once intent is translated into formal constraints, the placement stage can be formulated as optimization. Given infrastructure $\mathcal{G} = (\mathcal{N}, \mathcal{E})$ with $m$ nodes and services $\mathcal{S} = \{s_1, \ldots, s_{|\mathcal{S}|}\}$ with resource demands $\mathbf{d}_j = (d_j^c, d_j^m) \in \mathbb{R}_+^2$, placement $\pi: \mathcal{S} \to \mathcal{N}$ minimizes
\begin{equation}
\pi^*(t) = \arg\min_\pi \left[ J_{qos}(\pi, t; \mathsf{req}(t)) + \lambda R(\pi, t) + \mu E(\pi, t) \right],
\end{equation}
subject to capacity constraints $\sum_{j: \pi(s_j) = n_i} \mathbf{d}_j / \mathbf{r}_i \leq \mathbf{1}$ where $\mathbf{r}_i$ is node capacity. Load balancing cost $R(\pi, t) = \sum_{i=1}^m \|\mathbf{u}_i(\pi, t)\|_2^2$ penalizes imbalance, energy cost $E(\pi, t) = \sum_{i=1}^m p_i(\pi, t)$ accounts for power, and $\lambda, \mu \geq 0$ trade off objectives.

Even this single stage presents significant challenges. The decision space grows combinatorially ($|\mathcal{N}|^{|\mathcal{S}|}$ placements), objectives conflict dynamically, and optimal policies depend on workload characteristics varying across deployments. \acp{LLM} offer an alternative through intent-driven orchestration, where natural language processing enables Stage~(i), zero-shot reasoning addresses Stage~(ii), and persistent agents enable Stage~(iii)~\cite{kim2025kubeteus, netllm2024}.

\subsubsection{LLM-Based Reasoning with Structured Outputs}

\acp{LLM} map natural language intent $\mathsf{req}(t)$ to orchestration decisions through transformer-based reasoning. These models employ the same attention mechanism introduced in Section~\ref{subsec:ch2_attention_mechanism}, but applied to text sequences rather than time-series data. Given prompt $\mathbf{p}(t)$ encoding intent, infrastructure state $\mathbf{u}(t)$, and constraints, the \ac{LLM} generates output
\begin{equation}
\mathbf{y}(t) = \text{LLM}(\mathbf{p}(t)),
\end{equation}
where $\mathbf{y}(t)$ contains service specifications and placement decisions.

Unstructured \ac{LLM} outputs risk hallucinations that violate physical constraints. Structured outputs with validation schemas enforce feasibility. Given infrastructure state $\mathbf{u}(t) = [\mathbf{u}_1(t), \ldots, \mathbf{u}_m(t)]$, a validation schema for agent $i$ is generated as
\begin{equation}
\mathcal{V}_i(t) = \text{Schema}(\mathbf{u}(t), \text{role}_i),
\end{equation}
where $\text{role}_i$ specifies the agent's designated role. Validation of structured output $\mathbf{z}$ is defined as
\begin{equation}
\text{Valid}(\mathbf{z}, \mathcal{V}_i(t)) = \begin{cases}
1, & \text{if } \mathbf{z} \in \mathcal{V}_i(t) \\
0, & \text{otherwise}
\end{cases},
\end{equation}
where $\mathbf{z}$ is accepted only if it conforms to the dynamically generated constraints. The schema $\mathcal{V}_i(t)$ adapts dynamically, so when $\mathbf{u}_i(t) + \mathbf{d}_j/\mathbf{r}_i > \mathbf{1}$, node $n_i$ is excluded from valid placements, preventing capacity violations. Schema-validated outputs address single decisions, but production orchestration requires continuous monitoring and iterative refinement that single \ac{LLM} calls cannot provide.

\subsubsection{Autonomous Agents with Tool Access}

While \acp{LLM} with structured outputs enable single-step decision generation, real orchestration requires continuous monitoring, iterative refinement, and persistent context across interactions. Autonomous agents augment \acp{LLM}, adding capabilities that enable closed-loop control. The \ac{PARES} framework provides a capability taxonomy derived from analysis of existing agentic frameworks~\cite{sumers2024cognitive,Park2023GenerativeAgents,huang2023benchmarking,hao2023toolkengpt,zhang2023planning,promptreport2024}. For any autonomous agent $a_i$, \ac{PARES} defines five fundamental capabilities:
\begin{align}
\text{P}_i(\mathbf{o}_i(t)) &\rightarrow \text{Perceive environment state} \nonumber \\
\text{A}_i(\mathbf{T}_i) &\rightarrow \text{Act through tools} \nonumber \\
\text{R}_i(\mathbf{s}_i) &\rightarrow \text{Reason with internal state} \nonumber \\
\text{E}_i(\mathbf{m}_i) &\rightarrow \text{Evaluate outcomes} \nonumber \\
\text{S}_i(f_i) &\rightarrow \text{Sustain memory across interactions}. \label{eq:ch2_pares}
\end{align}
These five capabilities emerged as the common constructs distinguishing autonomous agents from simple scripts, namely observing environment state, executing actions through tools, reasoning about decisions, evaluating outcomes, and maintaining context across interactions.

The agent output is computed as
\begin{equation}
\mathbf{output}_i = f_i(\mathbf{o}_i(t), \mathbf{s}_i, \mathbf{m}_i, \mathbf{T}_i),
\end{equation}
where $f_i$ is the agent function, $\mathbf{o}_i(t)$ represents current observations, $\mathbf{s}_i$ is the internal reasoning state, $\mathbf{m}_i$ is historical memory, and $\mathbf{T}_i$ is the available tool set. Unlike stateless \ac{LLM} calls, agents maintain persistent state allowing closed-loop control and adaptive decision-making.

Generic agentic frameworks use trial-and-error learning through direct infrastructure experimentation~\cite{yao2023react}. For mission-critical orchestration, such exploratory approaches introduce significant reliability risks where failures cascade across distributed services. Beyond safety concerns, single-agent orchestration also creates sequential bottlenecks when handling complex multi-step workflows, motivating decomposition across multiple specialized agents.

\subsubsection{Multi-Agent Coordination}

Multi-agent systems address both scalability and safety limitations of single-agent orchestration by distributing decision-making across specialized agents. This reduces bottlenecks and cognitive overload while enabling domain-constrained reasoning. Formally, $\mathcal{A} = \{a_1, \ldots, a_q\}$ agents with roles $\mathcal{R} = \{\text{role}_1, \ldots, \text{role}_q\}$ where each $a_i$ has $\text{role}_i \in \mathcal{R}$ defining responsibilities.

Inter-agent communication follows message-passing, expressed as
\begin{equation}
\mathbf{msg}_{ij}(t) = g_{ij}(\mathbf{output}_i),
\end{equation}
where $g_{ij}$ transforms agent $a_i$'s output for agent $a_j$. This allows parallel processing, where intent parsing, state monitoring, and action planning occur simultaneously across specialized agents.

In orchestration specifically, multi-agent decomposition aligns with workflow phases. Intent agents parse natural language $\mathsf{req}(t)$ to structured goals, observe agents monitor $\mathbf{u}(t)$ for constraint violations, plan agents generate placement $\pi$ satisfying constraints, and act agents execute decisions through infrastructure APIs. Specialization maintains focused reasoning per domain, avoiding context-switching overhead. While specialization improves reasoning quality, the structured output mechanisms introduced earlier also provide quantifiable guarantees for reducing orchestration errors.

\subsubsection{Structured Output Error Reduction}

Beyond basic validation, structured outputs provide quantifiable error reduction. Consider three sets of possible actions. The set $\Omega_{all}$ contains all possible actions including impossible ones, $\Omega_{valid}(t)$ contains only feasible actions at time $t$, and $\Omega_{optimal}(t)$ contains the best feasible actions at time $t$, where $\Omega_{optimal}(t) \subseteq \Omega_{valid}(t) \subset \Omega_{all}$.

Without schema validation, an \ac{LLM} can choose any action from $\Omega_{all}$. The probability of making an error (choosing an infeasible action) is
\begin{equation}
P_{error}^{unconstrained} = 1 - \frac{|\Omega_{valid}(t)|}{|\Omega_{all}|}.
\label{eq:ch2_unconstrained_error}
\end{equation}
This probability is high because most random placements violate capacity constraints. With schema validation restricting choices to $\Omega_{valid}(t)$, the error probability becomes
\begin{equation}
P_{error}^{constrained} = 1 - \frac{|\Omega_{optimal}(t)|}{|\Omega_{valid}(t)|}.
\label{eq:ch2_constrained_error}
\end{equation}
The improvement $\Delta P_{error} = P_{error}^{unconstrained} - P_{error}^{constrained} > 0$ occurs because impossible actions are eliminated while retaining good actions, reducing orchestration failures. This error reduction applies to individual agent outputs, while the multi-agent decomposition described earlier provides additional benefits through distributed reasoning.

\subsubsection{Distributed Reasoning Benefits}

The benefits of distributing reasoning across specialized agents can be quantified mathematically. Research confirms that \acp{LLM} achieve better performance when complex reasoning is decomposed into intermediate steps rather than attempted monolithically~\cite{chainofthought2022}.

For mathematical analysis, let $p \in (0,1]$ be the baseline accuracy for focused single-step tasks, $\beta \in (0,1)$ the accuracy loss per step interaction, and $\eta$ the number of required sequential reasoning steps. When a single \ac{LLM} attempts all steps, accuracy degrades due to task complexity, yielding success probability
\begin{equation}
P_{monolithic} = p^\eta \cdot (1-\beta)^{\eta(\eta-1)}.
\label{eq:ch2_monolithic_success}
\end{equation}
In contrast, distributing reasoning between specialized agents maintains baseline accuracy $p$ for each designated step, giving
\begin{equation}
P_{distributed} = p^\eta.
\label{eq:ch2_distributed_success}
\end{equation}
The improvement factor $P_{distributed}/P_{monolithic} = 1/(1-\beta)^{\eta(\eta-1)}$ grows superexponentially with complexity $\eta$, providing mathematical justification for hierarchical multi-agent architectures.

Chapter~\ref{ch:distributed-agentic-ai} instantiates these concepts through \ac{AgentEdge}'s four-agent architecture, addressing distributed \ac{6G} edge-cloud orchestration through coordinated multi-agent workflows with simulation-based validation.


% Part II: State of the Art
\section{Part II: State of the Art}
\label{sec:ch2_part2}

Building on the foundational concepts established in Part I, this section reviews the state of the art in edge prediction, generalization approaches, and autonomous orchestration systems. Each review identifies specific gaps that motivate this dissertation's contributions.

% 2.4 Edge-Deployable Workload Prediction (State of the Art)
\subsection{Edge-Deployable Workload Prediction}
\label{sec:ch2_sota_aero}

Section~\ref{sec:ch2_ml_workload} introduced the machine learning architectures used for workload prediction and identified critical edge deployment constraints: model size below 1\ac{GB}, inference latency under 50ms, and sufficient accuracy for proactive orchestration. This section evaluates how existing approaches meet these requirements, reviewing: (i) recurrent and convolutional architectures for sequential modeling, (ii) graph-based and transformer approaches for complex dependencies, (iii) large language models and lightweight alternatives, and (iv) fundamental limitations that motivate the \ac{AERO} framework, detailed in Chapter~\ref{ch:adaptive-orchestration}.

\subsubsection{Recurrent and Convolutional Architectures}

Traditional time-series forecasting has relied on architectures designed to capture sequential patterns and local temporal features.

\textit{Recurrent Neural Networks:} \acp{RNN}~\cite{lai2018modelinglongshorttermtemporal, Jhin_2024, li2023learningpolarrepresentationextremeadaptive} have been important in time-series forecasting due to their ability to model sequential data. Recent works have proposed novel \ac{RNN} architectures and training strategies to tackle challenges like long-term dependencies and interpretability. Some models integrate autoregressive recurrent networks with probabilistic forecasting, while others combine \acp{CNN} and \acp{RNN} for improved multivariate prediction~\cite{li2024evogwp}.

\textit{Convolutional Neural Networks:} \acp{CNN}~\cite{wang2023micn} have been adapted for time-series forecasting by capturing local temporal patterns and hierarchical features. Research in this area focuses on specialized \ac{CNN} architectures for time-series data, multi-scale convolutional operations, and hybrid models combining \acp{CNN} with other techniques for improved performance. Approaches include deep architectures with residual links and fully-connected layers~\cite{donghao2024moderntcn}. However, \acp{CNN} struggle to capture long-range dependencies in time-series data due to their limited receptive field size~\cite{wu2023timesnettemporal2dvariationmodeling}. They also lack the ability to model sequential dependencies inherent in time-series, as they treat each input window independently. These sequential modeling limitations motivate the exploration of architectures that can capture more complex temporal and spatial dependencies.

\subsubsection{Graph-Based and Transformer Approaches}

Beyond recurrent and convolutional methods, advanced architectures use graph structures and attention mechanisms to capture complex temporal and spatial dependencies.

\textit{Graph Neural Networks:} For multivariate and spatio-temporal forecasting tasks, \acp{GNN}~\cite{yi2024fouriergnn, zhang2024irregular} model dependencies between different time-series as graphs. Building on this representation, recent studies combine temporal and spatial information to improve forecasting accuracy. These methods integrate \acp{GNN} with other deep learning techniques to capture complex dependencies. However, \acp{GNN} can be computationally expensive, especially for large-scale problems with complex dependency structures~\cite{jin2024survey}.

\textit{Transformer Architectures:} Transformer architectures~\cite{wu2023timesnettemporal2dvariationmodeling} have achieved success in time-series forecasting, particularly for long-term predictions~\cite{vaswani2017attention}. Recent works aim to improve the efficiency and effectiveness of Transformers by addressing temporal dependencies, capturing multi-scale patterns, and handling non-stationarity. Some models introduce methods to capture periodic patterns at different scales, while others focus on adapting to changing data distributions~\cite{jia2024witran}. However, Transformer models typically have a large number of parameters, often in the millions, leading to high computational requirements~\cite{liu2023nonstationary}. The computational demands of both \acp{GNN} and Transformers have prompted research into both larger foundation models and more efficient lightweight alternatives.

\subsubsection{Large Language Models and Lightweight Alternatives}

Recent trends explore both large-scale foundation models and efficient lightweight architectures for time-series forecasting.

\textit{\ac{LLM}-Based Forecasting:} An emerging trend uses \acp{LLM}~\cite{tang2023one, liu2023large} for time-series forecasting, taking advantage of their rich representations and transfer learning capabilities. However, due to the billions of parameters in \acp{LLM}, these approaches often need high computational resources~\cite{rasul2023forecasting}, making them less suitable for resource-constrained environments like edge-cloud systems.

\textit{\ac{MLP}-Based Models:} In contrast to parameter-heavy \acp{LLM}, \ac{MLP}-based models~\cite{zeng2023koopa, wu2023timemlp, liu2023fits} have gained attention due to their simplicity and computational efficiency. Recent research explores novel \ac{MLP} architectures that can capture temporal dynamics without complex attention mechanisms. Strategies include using mathematical theories and frequency domain representations to model temporal dependencies~\cite{zhou2023patchTST}. However, \ac{MLP} models such as \ac{SparseTSF}~\cite{sparseTSF} assume a fixed periodicity, limiting their ability to adapt to changing workload patterns. Both large-scale and lightweight approaches thus face distinct barriers to edge deployment that warrant systematic examination.

\subsubsection{Fundamental Limitations and Gaps}

The architectures reviewed above share fundamental limitations for edge deployment. Despite advances in \ac{RNN}, \ac{CNN}, \ac{GNN}, Transformer, \ac{LLM}, and \ac{MLP}-based forecasting, most approaches encounter critical constraints when deployed in resource-constrained environments~\cite{acmtimeseriesreview2024}.

The fundamental limitations across these architectures for edge deployment include: (i) high computational requirements that exceed edge device capabilities, (ii) inability to adapt to dynamic workload patterns without retraining, (iii) lack of specialized mechanisms for resource-constrained edge-cloud environments where low latency and energy efficiency are critical, and (iv) reliance on fixed periodicity assumptions in lightweight models (like \ac{SparseTSF}), which struggle to capture the variable patterns inherent in microservice workloads~\cite{10202641}. Addressing these constraints requires a different approach to model design.

\subsubsection{Addressing the Limitations with AERO}

\ac{AERO} directly addresses these limitations by rethinking model architecture for the edge. Instead of compressing general-purpose models, it introduces a purpose-built lightweight architecture with \textit{adaptive} period detection that dynamically identifies and exploits the inherent periodicity of microservice workloads, rather than assuming fixed periodic patterns. By combining adaptive period detection with sparse neural networks, \ac{AERO} achieves competitive accuracy while maintaining a computational footprint suitable for edge deployment. This enables the real-time, energy-efficient, and adaptive forecasting required for \ac{6G} edge-cloud orchestration, as detailed in Chapter~\ref{ch:adaptive-orchestration}.

While \ac{AERO} addresses edge deployment constraints for individual workloads, production environments require models that generalize across diverse application types without per-workload retraining. The following section examines this cross-application generalization challenge.



% 2.5 Cross-Application Generalization (State of the Art)
\subsection{Cross-Application Generalization}  
\label{sec:ch2_sota_omni}

Section~\ref{subsec:ch2_generalization_challenge} formalized the cross-dataset generalization gap and its operational consequences, while Section~\ref{subsec:ch2_attention_mechanism} introduced attention mechanisms as a potential solution for learning transferable temporal patterns. This section evaluates how existing approaches address these challenges, reviewing three categories of cross-application generalization approaches: (i) transfer learning and domain adaptation techniques, (ii) attention-based models, and (iii) specialized architectures and ensemble methods. Each category is analyzed for its fundamental limitations, concluding with how \ac{OmniFORE} addresses these gaps.

\subsubsection{Transfer Learning and Domain Adaptation}

Transfer learning approaches for cross-domain workload prediction include feature-based methods and domain-adaptive techniques that combine transfer learning with adversarial training.

\textit{Feature-Based Transfer Learning:} Transfer learning frameworks~\cite{wang2019transfer} use domain adaptation to learn shared and domain-specific features, enabling knowledge transfer across related distributions. These methods identify common patterns across domains while preserving domain-specific characteristics.

\textit{Domain-Adaptive Approaches:} Domain-adaptive methods combine transfer learning with adversarial training to extract domain-invariant features for effective zero-shot scenarios. These approaches learn representations that remain stable across different workload distributions.

However, these methods show critical limitations: (i) they struggle to determine which features to transfer when source and target distributions differ widely, (ii) they need extensive offline training with paired domain data, and (iii) they lack mechanisms for rapid adaptation to emerging workload types, limiting their applicability in dynamic edge-cloud environments. These fundamental constraints motivate exploring attention mechanisms as an alternative paradigm for addressing cross-application generalization.

\subsubsection{Attention-Based Models}

While transfer learning focuses on adapting features across domains, attention mechanisms offer an alternative approach by dynamically weighting temporal segments to capture complex dependencies across heterogeneous workload traces.

\textit{Attention-Enhanced Neural Networks:} These networks~\cite{zhang2022attention} integrate multi-head attention with recurrent layers to focus on critical patterns, selectively emphasizing informative temporal segments while suppressing noise.

\textit{Attention-Based Transfer Learning:} This approach~\cite{li2022robust} combines pre-trained feature extractors with domain-adaptive attention to highlight salient features, allowing models to identify transferable patterns across workload domains.

Despite these capabilities, attention-based approaches face three main drawbacks: (i) standard implementations rely on \ac{RNN} components with inherent "forgetting" properties for long sequences~\cite{zhou2021informer}, (ii) attention matrices scale quadratically with sequence length, creating significant computational bottlenecks for edge devices, and (iii) without proper guidance, attention may focus on random correlations rather than fundamental patterns that allow true generalization. These single-paradigm limitations have led researchers to explore whether combining multiple approaches might overcome the individual weaknesses of both transfer learning and attention mechanisms.

\subsubsection{Specialized Architectures and Ensemble Methods}

Beyond these single-paradigm methods, researchers have developed hybrid architectures and ensemble frameworks that combine multiple model types for improved generalization.

\textit{Hybrid \ac{CNN}-\ac{RNN} Architectures:} \ac{CNN}-\ac{RNN} hybrids~\cite{lai2018modelinglongshorttermtemporal} integrate convolutional layers for local dependency patterns with recurrent components for long-term temporal dependencies. These models use complementary strengths of different neural architectures.

\textit{Adaptive Ensemble Frameworks:} Adaptive ensemble frameworks~\cite{kim2024ensemble} dynamically adjust the contribution of multiple models based on real-time performance. These systems adapt model weights to changing workload characteristics.

\textit{Modern Convolutional Approaches:} Recent work has introduced ModernTCN~\cite{ModernTCN}, a pure convolution-based approach that updates traditional \acp{TCN} with larger kernel sizes and specialized convolution blocks. By decoupling temporal and feature processing, ModernTCN performs well across diverse time-series tasks while maintaining lower computational costs than attention-based methods. Its linear complexity makes it suitable for environments with limited resources. However, ModernTCN was designed for general time-series forecasting; its effectiveness on highly variable container workloads in edge-cloud networks has not been extensively evaluated in the original work.

\textit{Graph, Generative, and Emerging Methods:} More recent approaches incorporate Graph Neural Networks~\cite{li2024evogwp} for capturing service dependencies, \acp{GAN} \cite{esteban2017real} for data augmentation, and even Quantum Neural Networks~\cite{10531701}. While these methods explore novel computational paradigms, they continue to face limitations from their underlying \ac{RNN} and \ac{CNN} components, particularly in resource-constrained edge environments.

The fundamental limitations of these approaches include: (i) high computational requirements that exceed edge device capabilities, (ii) the treatment of workload traces as individual entities instead of finding generic patterns, and (iii) the overhead of maintaining multiple models or complex architectures that often outweighs generalization benefits in resource-constrained settings. These limitations create a fundamental scalability bottleneck: deploying generalized prediction typically requires training separate models for each microservice or workload type, resulting in an unsustainable "n-models for n-services" paradigm. The following subsubsection examines how \ac{OmniFORE} addresses these limitations through a different approach to workload pattern recognition.

\subsubsection{Addressing the Limitations with OmniFORE}

\ac{OmniFORE} overcomes the three generalization barriers identified above, namely feature transfer ambiguity, attention scalability, and architectural complexity. The approach moves beyond the feature-matching paradigm of transfer learning to identify fundamental behavioral patterns. By introducing temporal clustering as a precursor to attention mechanisms, it categorizes disparate workload traces into a finite set of universal behaviors. This allows the model to learn from representative archetypes rather than overfitting to individual traces, substantially mitigating the "n-models for n-services" scalability challenge. Chapter~\ref{ch:attention-driven-prediction} details how this pattern-based approach enables generalization across heterogeneous microservices without per-service retraining.

With edge-deployable prediction (\ac{AERO}) and cross-application generalization (\ac{OmniFORE}) established, the remaining challenge is translating these predictions into autonomous orchestration actions. The following section examines how existing systems attempt to bridge this gap between prediction and action.
% 2.6 Autonomous Edge-Cloud Orchestration (State of the Art)
\subsection{Autonomous Edge-Cloud Orchestration}
\label{sec:ch2_sota_agent}

Section~\ref{sec:ch2_intelligent_orchestration} established the theoretical foundations for \ac{AI}-driven orchestration, including the three-stage orchestration problem, the \ac{PARES} capability framework, and the mathematical benefits of multi-agent coordination. This section evaluates how existing systems implement these concepts, reviewing: (i) \ac{LLM}-based orchestration approaches for intent processing and code generation, (ii) agent-based orchestration systems spanning domain-specific applications, multi-agent coordination, and infrastructure management, and (iii) fundamental limitations that motivate \ac{AgentEdge}.

\subsubsection{LLM-based Orchestration}

\ac{LLM}-based approaches divide into two primary categories: intent processing and code generation.

\textit{Intent Processing and Translation:} Translating natural language intents into network configurations requires breaking down ambiguous requests into verifiable actions. Recent work~\cite{intent_driven_ran2021} uses few-shot learning to break down and validate intent in next-generation networks. Other studies~\cite{IntentBasedNetworkManagementDRCN2024} show intent extraction capabilities for \ac{6G} core networks. Advanced approaches use a four-stage methodology to convert natural language intentions into Network Service Descriptors using Code Llama \ac{LLM}~\cite{LLMIntentServiceConfig2024}, incorporating few-shot learning, similarity-based retrieval, closed-loop validation, and human feedback for continuous improvement.

\textit{Code Generation and Configuration Management:} Beyond intent extraction, automating network configuration requires generating correct code while addressing privacy and error-handling concerns. Recent studies~\cite{EnhancingNetworkManagementMSR2023} present task-specific code generation for network graph analysis and manipulation, combining application-specific prompts with program synthesis techniques to address explainability, scalability, and privacy challenges. Complementary work~\cite{IntentNetworksLLMFramework2024} introduces a modular framework that addresses data privacy and \ac{LLM} errors in router and firewall configurations.

However, current \ac{LLM}-based orchestration systems face critical limitations~\cite{NEURIPS2024_6aebba00, struggle_with_real_time_adaptation2025}: model hallucinations that generate incorrect configurations, inference latencies that prevent real-time adaptation, insufficient tool integration for infrastructure \acp{API}, and monolithic architectures unsuitable for distributed edge-cloud environments. The monolithic nature of these systems creates a fundamental mismatch with the distributed requirements of edge-cloud orchestration.

\subsubsection{Agent-based Orchestration}

Building upon these \ac{LLM} foundations, agent-based orchestration systems represent the next evolution, using autonomous \ac{AI} agents for coordinated service orchestration across heterogeneous environments.

\textit{Domain-Specific Agent Approaches:} Implementations target specialized networking applications with customized agent architectures. In wireless network management, researchers~\cite{WirelessAgentLLM2025} demonstrate intelligent network management through \ac{LLM}-powered agents. For fault management, other work~\cite{LLMAgentsProactiveFaultTolerant2024} introduces an \ac{LLM}-based agent system with integration of memory, knowledge base, and network infrastructure. This system autonomously manages the complete fault management cycle through proactive traffic monitoring, delay detection, and flow rerouting. For network planning applications, researchers~\cite{RadioMapAgentsNetworkPlanning2025} present a three-module framework with specialized tools for geographic map retrieval, environment creation, and automated radio map generation. However, these approaches rely on single-agent architectures that cannot distribute reasoning across specialized roles, lack formal capability definitions, and execute actions directly without pre-validation in simulated environments.

\textit{Agent Coordination and Collaboration:} Moving beyond single-agent approaches, researchers have developed frameworks for multi-agent cooperation and coordinated decision-making. Recent work~\cite{HermesLLMFramework2024} uses a chain-of-agent approach that separates network modeling into specialized roles using \ac{YAML}-based blueprints with mathematical formulas and pseudocode. Other approaches~\cite{MaestroLLMCollaborative2024} present multi-agent systems with \ac{LLM}-based agents for conflict resolution and resource negotiation on the business, service, and network planes. Additional work~\cite{xiao2025agentnet} uses agents across the application, physical, and network layers, evaluated through industrial automation and virtual reality applications. These frameworks typically use flat coordination structures without hierarchical organization, accept \ac{LLM} outputs without constraint verification, and test plans through trial-and-error on production infrastructure rather than simulation.

\textit{Infrastructure Orchestration and Service Management:} Finally, infrastructure orchestration approaches focus on agentic implementations for distributed service management and deployment. Recent work~\cite{DataConnectorAgentFramework2024} uses machine learning models for \ac{QoS} predictions and allows autonomous migration of containerized applications in Kubernetes-based edge-cloud environments. Other work~\cite{service_orch_2025} presents an intent-based orchestration approach using four specialized agents with distributed coordination. This system employs Model-Context Protocol integration for tool standardization and uses fine-tuned Mistral-7B and Qwen3-0.6B models through \ac{QLoRA}~\cite{dettmers2023qlora} for Kubernetes configuration tasks. These systems assume cloud-centric architectures with reliable high-bandwidth connectivity, limiting applicability to heterogeneous edge-cloud environments, and address deployment but not migration, scaling, or power management. Despite their varied approaches, these agent-based systems share underlying assumptions that constrain their applicability to distributed edge-cloud scenarios.

\subsubsection{Fundamental Limitations and Research Gaps}

The \ac{LLM}-based and agent-based approaches reviewed above share critical gaps for distributed edge-cloud environments. They focus on centralized cloud infrastructures, target narrow applications without addressing complete service lifecycle management, and lack formal agent capability definitions and coordination mechanisms for distributed multi-agent collaboration.

Beyond domain-specific systems, generic agentic approaches from the broader \ac{ML} literature~\cite{yao2023react, LATS} such as \ac{LATS} employ trial-and-error learning requiring direct experimentation on production infrastructure. This creates reliability concerns for deployments where failures cascade across distributed services. In contrast, traditional optimization methods offer deterministic guarantees but have the converse limitation: they solve placement problems given formally specified objectives, yet cannot interpret natural language intents or handle partial infrastructure states~\cite{6g_struggles_orchestration2025}. These gaps reveal a persistent safety-autonomy trade-off: autonomous agents offer adaptability but risk production failures, while rule-based systems provide safety but lack flexibility.

\subsubsection{Addressing the Limitations with AgentEdge}

\ac{AgentEdge} addresses this safety-autonomy trade-off through its \ac{ActSimCrit} methodology, where agents Act in simulation, observe outcomes through Simulation, and receive feedback through Critique before production deployment. Instead of choosing between unsafe exploration (\ac{RL}/agents) or rigid heuristics (rule-based), \ac{AgentEdge} introduces a third path: validated autonomy. By integrating a high-fidelity digital twin directly into the agent's reasoning loop, it enables risk-free exploration of complex orchestration strategies. This allows agents to learn from failure in simulation rather than production, combining the adaptability of generative AI with improved reliability through pre-deployment validation. Chapter~\ref{ch:distributed-agentic-ai} presents this framework, demonstrating how distributed specialized agents can safely orchestrate mission-critical \ac{6G} infrastructure.


% Summary
\section{Summary: From Gaps to Contributions}
\label{sec:ch2_summary}

The foundational concepts and state-of-the-art analysis presented in this chapter reveal three distinct research gaps that prevent autonomous edge-cloud orchestration.

First, existing prediction models achieve high accuracy through large architectures unsuitable for edge deployment, leaving a gap between prediction capability and edge constraints (Section~\ref{sec:ch2_sota_aero}). Second, current approaches require per-application training, creating unsustainable operational overhead when managing heterogeneous microservices at scale (Section~\ref{sec:ch2_sota_omni}). Third, autonomous orchestration systems either lack safety guarantees or sacrifice flexibility for reliability, leaving the safety-autonomy paradox unresolved (Section~\ref{sec:ch2_sota_agent}).

The following three chapters address these gaps directly. Chapter~\ref{ch:adaptive-orchestration} introduces \ac{AERO}, a model with fewer than 1,000 parameters that achieves edge-deployable prediction through adaptive period detection. Chapter~\ref{ch:attention-driven-prediction} presents \ac{OmniFORE}, which enables cross-application generalization through temporal clustering and attention mechanisms. Chapter~\ref{ch:distributed-agentic-ai} develops \ac{AgentEdge}, a multi-agent framework that achieves validated autonomy through simulation-first decision-making. Together, these contributions form a complementary pipeline from lightweight prediction through scalable generalization to autonomous orchestration.
